% Candidate fragment generated from 06_执行台账/cross_version_20261012/.
% Verify figure filenames and section numbering against the Overleaf main version.
\subsection{Conditioned cross-version comparison}
The baseline, an independently loaded BBRv1 module, BBR2, and BBRv3 runs were analyzed with the same condition
keys, full-run receiver-summary window, and five-run repetition rule.
Although configuration files declare a 30-s warmup, the archived primary
analysis uses the full 180-s transfer. A separate interval-based
sensitivity analysis retains the requested post-warmup window. The Ubuntu in-tree
baseline is retained as an implementation-specific reference. BBRv1,
BBR2, and BBRv3 are identified by their source/kernel markers; their runtime
registration names are \texttt{bbr1}, \texttt{bbr2}, and \texttt{bbr},
respectively. The BBRv1 module uses the recorded 6.13 compatibility patch
for the unexported TSO helper.

Across the tested 100-Mbps matrix, equal-RTT total goodput remains close to
the bottleneck for all four implementations, while flow shares and Jain
fairness vary with RTT asymmetry, queue setting, and implementation. Under
the 40- and 80-ms second-flow conditions, the observed long-RTT share differs across implementations. These results support a
conditional comparison over the tested matrix; they do not establish a
universal version ranking or identify unmeasured internal state.

For the no-ECN BBR/CUBIC comparison, total goodput is similar across the
four implementation environments, while the BBR/CUBIC share and Jain
index remain condition dependent. The corresponding values and raw inputs
are archived in \texttt{cross\_version\_comparison.csv}; figures are generated by
\texttt{analyze\_cross\_version.py}.
